\documentclass[journal]{IEEEtran}

\usepackage{cite}
\ifCLASSINFOpdf
  \usepackage[pdftex]{graphicx}
\else
  \usepackage[dvips]{graphicx}
\fi
\usepackage{amsmath}
\usepackage{amssymb}
\usepackage{array}
\usepackage{booktabs}
\hyphenation{op-tical net-works semi-conduc-tor}

\begin{document}

\title{Dynamic Spatial Stabilization and Trajectory Tracking of a Spherical Agent on a 6-DOF Robotic Platform}

\author{Phillipp Gery (pgery), Kevin Biju Mathew (kb)%
  \thanks{College of Engineering, Purdue University, Indiana, USA.}}

\markboth{Course Project Milestone 1: Robotics Control, March~2026}%
{Gery \MakeLowercase{\textit{et al.}}: Autonomous Marble Balancing}

\maketitle

\begin{abstract}
  This report presents the implementation of a real-time discrete-time Linear Quadratic Regulator (LQR)
  controller for stabilizing a marble on the end-effector plate of a UR5e 6-DOF robotic arm in simulation.
  The simulation environment is built on Gazebo Classic, ROS2 Humble, and MoveIt2 Servo. An 8-dimensional
  state-space model captures the coupled ball-plate dynamics and the robot's first-order (PT1) actuator lag.
  The LQR gain is computed offline via the Discrete Algebraic Riccati Equation and applied at 30\,Hz.
  Preliminary results demonstrate successful marble stabilization, autonomous recovery after plate departure,
  and Lissajous trajectory tracking. Future work includes a Soft Actor-Critic residual policy, feedforward
  tilt compensation, and hardware transfer to a physical UR5e arm.
\end{abstract}

\begin{IEEEkeywords}
  ROS 2, Gazebo, LQR, MoveIt Servo, UR5e, Trajectory Tracking, Ball-Plate System.
\end{IEEEkeywords}

% =============================================================================
\section{Introduction}
% =============================================================================

\IEEEPARstart{P}{recision} stabilization of unconstrained objects is a fundamental challenge in robotics.
This project implements real-time feedback control of a UR5e 6-DOF robotic arm that balances a marble on a
flat plate rigidly attached to the robot's end-effector. The arm tilts the plate via small angular velocity
commands, keeping the marble centered. The full system operates in simulation using
\textbf{Gazebo Classic}, \textbf{ROS2 Humble}, and \textbf{MoveIt2 Servo}.

The problem couples two dynamical systems: the rolling dynamics of the marble on the plate, and the servo
dynamics of the arm responding to velocity commands. Marble acceleration depends on plate tilt angle, while
the arm's response to commands is subject to inherent actuator lag. A classical discrete-time
\textbf{Linear Quadratic Regulator (LQR)} serves as the primary controller, designed around a physics model
that explicitly captures both the ball-plate coupling and the robot's first-order (PT1) delay.

Unlike traditional ``Labyrinth'' setups where the board orientation is directly actuated, this project
employs \textit{active surface control}: a 6-DOF arm adjusts end-effector orientation in real time while
simultaneously being capable of tracing arbitrary 3D paths with the TCP. Optional modes extend the base
system to trajectory tracking: the marble can follow a Lissajous curve setpoint, the robot TCP can trace a
figure-eight path, and a Soft Actor-Critic (SAC) residual controller can augment the LQR output. This
report covers the simulation setup and LQR classical controller constituting Milestone 1.

% =============================================================================
\section{Methods}
% =============================================================================

\subsection{Simulation Environment}

\subsubsection{Robot and World}
The UR5e arm is loaded into Gazebo Classic via a URDF xacro that uses the standard \texttt{ur\_description}
macros. A $0.40\,\mathrm{m} \times 0.40\,\mathrm{m} \times 5\,\mathrm{mm}$ acrylic plate is rigidly
attached to the robot's \texttt{tool0} link. A virtual frame \texttt{plate\_tcp} sits at the plate's top
surface and serves as the control reference frame. The simulated world contains only the robot, the plate,
and the marble.

\subsubsection{Marble}
The marble is a $15\,\mathrm{mm}$ radius, $50\,\mathrm{g}$ sphere defined in SDF format. Gazebo's ODE
physics engine simulates rolling contact on the plate with a friction coefficient of 0.6. The marble is
dynamically spawned above the plate after the arm reaches its home pose via the \texttt{/spawn\_entity}
Gazebo service. The spawn position is computed at runtime from a TF lookup of the \texttt{plate\_tcp}
frame, so no coordinates are hardcoded.

\subsubsection{State Sensing}
Three sensing channels provide the full 8-D state:

\begin{itemize}
  \item \textbf{Marble position and velocity:} The \texttt{libgazebo\_ros\_p3d} Gazebo plugin publishes
    ground-truth odometry at 50\,Hz on \texttt{/marble/odom}. Velocity is estimated by numerically
    differentiating successive positions and filtering with an exponential moving average
    ($\tau = 0.08\,\mathrm{s}$).
  \item \textbf{Plate orientation:} Obtained by querying the TF tree for the
    \texttt{world} $\rightarrow$ \texttt{plate\_tcp} transform and converting the quaternion to roll/pitch
    Euler angles at the odometry rate ($\approx$50\,Hz).
  \item \textbf{Plate angular velocity:} Computed by numerically differentiating the TF-derived Euler
    angles ($\alpha$, $\beta$) in the same 50\,Hz odometry callback and smoothed with the same exponential
    moving average. An earlier Jacobian-based approach ($\mathbf{J}_\mathrm{rot}\,\dot{\mathbf{q}}$ from
    100\,Hz joint states) was abandoned because it produced a sign inconsistency relative to the TF-derived
    angles at the robot's home configuration, causing oscillation in the pitch channel.
\end{itemize}

\subsubsection{Initialization Sequence}
The system uses an event-driven launch sequence (no fixed time delays). After Gazebo and MoveIt load, a
\texttt{go\_to\_pose} node commands the arm to its home Cartesian pose via inverse kinematics and a
\texttt{JointTrajectoryController} action. When that node exits, the \texttt{marble\_spawner} fires, drops
the marble, and exits — triggering the controller and auxiliary nodes via ROS2 \texttt{OnProcessExit}
launch event handlers.

% -----------------------------------------------------------------------------
\subsection{Classical Controller: Discrete-Time LQR}
% -----------------------------------------------------------------------------

\subsubsection{Physical Model}
The ball-plate system is modeled as a planar coupled system. Under the small-angle rolling assumption for a
solid sphere, ball acceleration along each axis is proportional to plate tilt:
%
\begin{equation}
  a_\mathrm{ball} = -C\,\theta_\mathrm{plate}, \qquad
  C = \frac{Mg}{M + I/r^2} \approx 7.0\,\mathrm{m/s^2},
  \label{eq:ball_accel}
\end{equation}
%
where $M = 0.05\,\mathrm{kg}$, $I = \tfrac{2}{5}Mr^2$, $r = 0.015\,\mathrm{m}$, and
$g = 9.81\,\mathrm{m/s^2}$.

The robot arm's response to angular velocity commands is modeled as a first-order (PT1) lag with time
constant $T_\mathrm{robot} = 0.35\,\mathrm{s}$:
%
\begin{equation}
  \dot{\omega} = -\frac{\omega - \omega_\mathrm{cmd}}{T_\mathrm{robot}}.
  \label{eq:pt1}
\end{equation}
%
This captures the combined effect of communication latency, controller bandwidth, and joint servo dynamics.

\subsubsection{State Vector (8-D)}
The full state vector is:
%
\begin{equation}
  \mathbf{x} = \begin{bmatrix}
    x & v_x & y & v_y & \alpha & \omega_\alpha & \beta & \omega_\beta
  \end{bmatrix}^\top,
  \label{eq:state}
\end{equation}
%
where $(x,y)$ are marble positions on the plate, $(v_x, v_y)$ are marble velocities, $\alpha$ (pitch)
and $\beta$ (roll) are plate tilt angles governing $x$- and $y$-ball motion respectively, and
$\omega_\alpha$, $\omega_\beta$ are their rates of change (PT1 states). The two axes are decoupled in the
model, forming an $8\times 8$ system.

\subsubsection{Continuous-Time System Matrices}
The continuous-time state-space matrices $\mathbf{A} \in \mathbb{R}^{8\times 8}$ and
$\mathbf{B} \in \mathbb{R}^{8\times 2}$ are:
%
\begin{equation}
  \mathbf{A} = \begin{bmatrix}
    0 & 1 & 0 & 0 &  0    & 0 &  0    & 0 \\
    0 & 0 & 0 & 0 & -C    & 0 &  0    & 0 \\
    0 & 0 & 0 & 1 &  0    & 0 &  0    & 0 \\
    0 & 0 & 0 & 0 &  0    & 0 & -C    & 0 \\
    0 & 0 & 0 & 0 &  0    & 1 &  0    & 0 \\
    0 & 0 & 0 & 0 & -1/T  & 0 &  0    & 0 \\
    0 & 0 & 0 & 0 &  0    & 0 &  0    & 1 \\
    0 & 0 & 0 & 0 &  0    & 0 & -1/T  & 0
  \end{bmatrix},
  \label{eq:A}
\end{equation}
%
\begin{equation}
  \mathbf{B} = \begin{bmatrix}
    0 & 0 \\ 0 & 0 \\ 0 & 0 \\ 0 & 0 \\
    0 & 0 \\ 1/T & 0 \\ 0 & 0 \\ 0 & 1/T
  \end{bmatrix},
  \label{eq:B}
\end{equation}
%
where $T = T_\mathrm{robot}$. Control inputs
$\mathbf{u} = [\omega_{\alpha,\mathrm{cmd}},\; \omega_{\beta,\mathrm{cmd}}]^\top$ are angular velocity
commands sent to MoveIt Servo.

\subsubsection{Discretization (Zero-Order Hold)}
The continuous system is discretized using the matrix-exponential method at $\Delta t = 1/30\,\mathrm{s}$:
%
\begin{equation}
  \begin{bmatrix} \mathbf{A}_d & \mathbf{B}_d \\ \mathbf{0} & \mathbf{0} \end{bmatrix}
  = \exp\!\left(
    \begin{bmatrix} \mathbf{A} & \mathbf{B} \\ \mathbf{0} & \mathbf{0} \end{bmatrix} \Delta t
  \right),
  \label{eq:zoh}
\end{equation}
%
implemented via \texttt{scipy.linalg.expm}, ensuring accurate discretization of the coupled dynamics.

\subsubsection{LQR Gain Computation}
The discrete-time LQR gain $\mathbf{K}$ is found by solving the Discrete Algebraic Riccati Equation
(DARE):
%
\begin{align}
  \mathbf{P} &= \texttt{solve\_discrete\_are}(\mathbf{A}_d, \mathbf{B}_d, \mathbf{Q}, \mathbf{R}), \\
  \mathbf{K} &= (\mathbf{R} + \mathbf{B}_d^\top \mathbf{P}\, \mathbf{B}_d)^{-1}
                 \mathbf{B}_d^\top \mathbf{P}\, \mathbf{A}_d.
  \label{eq:dare}
\end{align}
%
The gain $\mathbf{K} \in \mathbb{R}^{2\times 8}$ is computed once at node startup.

\textbf{Cost matrices:}
%
\begin{align}
  \mathbf{Q} &= \mathrm{diag}(100,\;100,\;200,\;400,\;5,\;0.5,\;5,\;1), \label{eq:Q} \\
  \mathbf{R} &= 5.0\,\mathbf{I}_{2\times 2}. \label{eq:R}
\end{align}
%
Position weights ($Q_0, Q_2$) penalize centering error. Velocity weights ($Q_1, Q_3$) provide damping;
$Q_3 = 400$ is elevated relative to $Q_1$ because the Y axis is driven at twice the frequency in TCP
Lissajous mode, requiring stronger damping. Angle and rate weights ($Q_4$--$Q_7$) are kept small to allow
aggressive plate tilting without excessive control effort cost.

\subsubsection{Online Control Loop (30\,Hz)}
At each control tick the node:
\begin{enumerate}
  \item Assembles the 8-D state from filtered odometry and TF data.
  \item Computes the state error $\mathbf{e} = \mathbf{x} - \mathbf{x}_\mathrm{des}$ (setpoint is zero by
    default, or a Lissajous point).
  \item Applies LQR: $\mathbf{u} = -\mathbf{K}\,\mathbf{e}$.
  \item Saturates: $\mathbf{u} = \mathrm{clip}(\mathbf{u},\,-45^\circ/\mathrm{s},\,+45^\circ/\mathrm{s})$.
  \item Publishes as a \texttt{TwistStamped} angular velocity command to MoveIt Servo via the mux.
\end{enumerate}
MoveIt Servo converts the Cartesian angular velocity command into joint velocity commands at 30\,Hz,
which are then tracked by the UR5e's \texttt{JointTrajectoryController}.

\subsubsection{Mux Controller}
A priority multiplexer sits between the LQR controller and Servo. Manual operator commands override LQR
output, reverting to automatic control after 0.5\,s of inactivity. This allows safe operator intervention
without restarting the system.

% =============================================================================
\section{Preliminary Experiment Results}
% =============================================================================

The following describes observed system behavior during initial simulation runs.

\textbf{Homing and spawning:} The arm reliably reaches the home pose
(TCP at approximately $[0.233,\;0.25,\;0.8]\,\mathrm{m}$ in the world frame) via inverse kinematics before
the marble is dropped. The TF-based spawn calculation correctly places the marble 80\,mm above the plate
surface regardless of minor home pose variation. Landing detection (5 consecutive odometry ticks within
$\pm3\,\mathrm{cm}$ of the plate surface) triggers Servo activation without false positives.

\textbf{Stabilization — baseline LQR:} With default Q/R weights and no trajectory tracking, the LQR
controller successfully stabilizes the marble near the plate center. The marble settles from an initial drop
within approximately 2--3 seconds, with residual oscillation on the order of $\pm1$--$2\,\mathrm{cm}$ in
steady state. Both X and Y axes stabilize, though Y shows slightly more oscillation in early runs.

\textbf{Y-axis oscillation:} When the TCP Lissajous mode is active (the robot TCP traces a figure-eight at
amplitude 0.30\,m, period 12\,s, with Y at $2\times$ frequency), the increased plate acceleration in the Y
direction drives the marble further off center. Raising $Q_3$ (marble Y velocity cost) from 100 to 400
reduced this oscillation significantly. Raising $Q_2$ (marble Y position cost) to 200 further improved
centering. These changes are reflected in the current default weights.

\textbf{Lissajous trajectory tracking:} With marble Lissajous setpoints enabled, the controller
successfully causes the marble to follow a figure-eight pattern on the plate. Tracking error increases with
setpoint amplitude and frequency, as expected from the linear model's neglect of large-angle nonlinearities.

\textbf{Fell-off recovery:} When the marble rolls off the plate edge, the system correctly detects the fall
(marble $z$ drops more than 5\,cm below the plate surface), stops Servo, re-homes the arm, and re-spawns
the marble autonomously. This cycle completes in approximately 8--10 seconds.

% =============================================================================
\section{Future Work}
% =============================================================================

\subsection{Reinforcement Learning Residual Controller}
A Soft Actor-Critic (SAC) residual policy is partially implemented. The training environment
(\texttt{ball\_plate\_env.py}) simulates the same PT1 ball-plate physics with Coulomb/viscous friction and
domain randomization (mass, friction, time constant). A curriculum trains the agent through four stages of
increasing difficulty (residual authority $\pm2 \rightarrow \pm20\,^\circ/\mathrm{s}$, with Lissajous and
perturbations added progressively). The next step is to complete training, evaluate transfer to the Gazebo
simulation, and tune the LQR + RL blending.

\subsection{Feedforward Tilt Compensation}
The TCP Lissajous node computes and publishes a feedforward tilt vector (\texttt{/tcp/lissajous\_ff\_tilt})
intended to pre-compensate the pseudo-force the marble experiences during TCP acceleration. The feedforward
gain is currently disabled. The correct sign and magnitude need to be verified experimentally by comparing
marble trajectory error with and without compensation.

\subsection{Nonlinear / Adaptive Control}
The linear model assumes small plate angles and ignores friction. At higher TCP velocities or Lissajous
amplitudes these nonlinearities become significant. A planned extension is either to add a friction
compensation term using identified static and kinetic friction coefficients, or to switch to a nonlinear
model predictive controller (NMPC) that handles large-angle kinematics.

\subsection{Hardware Transfer}
The simulation is designed to match real UR5e hardware: the PT1 delay model uses a measured actuator time
constant, and \texttt{libgazebo\_ros\_p3d} is replaceable with a physical vision system (e.g., an overhead
camera with blob detection). The next milestone will involve validating the controller on a physical UR5e
arm, assessing whether the PT1 model adequately captures real actuator dynamics and whether the LQR gains
transfer without retuning.

\subsection{Kalman Filter for State Estimation}
Marble velocity and plate angular velocity are currently estimated by numerically differentiating position
measurements filtered with a simple exponential moving average. A planned improvement replaces this with a
\textbf{Kalman Filter (KF)} or \textbf{Extended Kalman Filter (EKF)} that fuses odometry, joint state, and
TF data using the ball-plate PT1 model as the process model. This will provide optimal state estimates under
Gaussian noise assumptions, reduce derivative noise without adding excessive lag, and improve controller
performance — particularly for velocity terms that currently require high Q weights to compensate for
estimation error.

\subsection{Quantitative Performance Metrics}
Current evaluation is qualitative. Future work will compute: RMS centering error, mean time to stabilize
from drop, maximum stable Lissajous amplitude, and a comparative analysis of LQR-only versus LQR+RL
performance on identical disturbance sequences.

% =============================================================================
\end{document}
